% Part: second-order-logic
% Chapter: sol-and-sets
% Section: introduction

\documentclass[../../../include/open-logic-section]{subfiles}

\begin{document}

\olfileid[hi]{sol}{set}{int}

\olsection{परिचय}

चूँकि द्वितीय-कोटि तर्कशास्त्र में फलनों के साथ-साथ प्रांत के उपसमुच्चयों
पर भी परिमाणीकरण किया जा सकता है, इसलिए यह अपेक्षित है कि कम-से-कम कुछ
समुच्चय-सिद्धांत द्वितीय-कोटि तर्कशास्त्र में विकसित किया जा सके। ``विकसित
करने'' से हमारा आशय है कि समुच्चय-सैद्धांतिक गुणों और कथनों को
द्वितीय-कोटि तर्कशास्त्र में व्यक्त किया जा सकता है, और यह समुच्चयों के लिए
किसी विशेष अतार्किक शब्दावली (जैसे समुच्चय-सिद्धांत का सदस्यता
!!{predicate}) के बिना संभव है। उदाहरणार्थ, हम समुच्चयों के संघ और
सर्वनिष्ठ तथा उपसमुच्चय-संबंध को परिभाषित कर सकते हैं; साथ ही समुच्चयों के
आकारों की तुलना और कैंटर के प्रमेय जैसे परिणामों का कथन भी कर सकते हैं।

\end{document}
